\documentclass[12pt]{article}
\usepackage[margin=1in]{geometry}
\usepackage{amsmath,amssymb,graphicx,xcolor,url,hyperref,breakurl,enumerate}


\title{Project Ideas: UConn--PCS: Data Science, 2023}
\begin{document}
\maketitle

\section{Probabilities with Simulation}
You may choose to use simulation to find answers to one (or more) of the
following problems.

\begin{enumerate}[1.]
\item A King offers a death row prisoner a chance to live by playing a simple
  game. He gives 50 black marbles, 50 white marbles, and 2 empty bowls to the
  prisoner. The King asks the prisoner to divide these 100 marbles into these 2
  bowls in any way as long as all the marbles are used. He will then blindfold
  the prisoner and mix the bowls around. After that, the prisoner can choose one
  bowl and pick one marble randomly.  If the marble is white the prisoner will
  live, but if the marble is black the prisoner will die. Use simulation to find
  the probability of living for the following ways of dividing.
  \begin{enumerate}
  \item Put the 50 white balls in one bowl and the 50 black balls in another
    bowl.
  \item Put one white ball in one bowl and the other 99 bowl balls in another
    bowl.
  \end{enumerate}
  
\item Tom and Jerry, each put 30 dollar in a jackpot to start a game. Suppose
  that they have equal chance to win and the one who wins three times takes all
  the 60 dollars. Now, Tom has won twice and Jerry has won once, but then
  something happens and the game must be stopped. How should they split the 60
  dollars according their probabilities of winning if the game was finished.

\item Three players enter a room and a red or blue hat is placed on each
  person's head.  The color of each hat is determined by an independent coin
  toss. No communication of any sort is allowed, except for an initial strategy
  session before the game begins.  Once they have had a chance to look at the
  other hats but not their own, the players must simultaneously guess the
  color of their own hats or pass. The puzzle is to find a group strategy that
  maximizes the probability that at least one person guesses correctly and
  no-one guesses incorrectly.
  \begin{enumerate}
  \item One strategy is for the group to agree that one person should guess and
    the others pass. What is the probability of success?
  \item Here is another strategy. For each person, guess the color he/she does
    not see if he/she see same colors, and pass is he/she sees different
    colors. What is the probability of success?
  \end{enumerate}
  
\item Henry has been caught stealing cattle, and is brought into town for
  justice. The judge is his ex-wife Gretchen, who wants to show him some
  sympathy, but the law clearly calls for two shots to be taken at Henry from
  close range. To make things a little better for Henry, Gretchen tells him she
  will place two bullets into a six-chambered revolver in successive order. She
  will spin the chamber, close it, and take one shot. If Henry is still alive,
  she will then either take another shot, or spin the chamber again before
  shooting. Henry is a bit incredulous that his own ex-wife would carry out the
  punishment, and a bit sad that she was always such a rule follower. He steels
  himself as Gretchen loads the chambers, spins the revolver, and pulls the
  trigger. Whew! It was blank. Then Gretchen asks, ``Do you want me to pull the
  trigger again, or should I spin the chamber a second time before pulling the
  trigger?'' What should Henry choose?
  
\item In a prison, there are 100 death row prisoners who are numbered from 1 to
  100, and there is a room with 100 drawers labeled from 1 to 100. The director
  randomly puts one prisoner's number in each closed drawer and offers a last
  chance. The prisoners enter the room, one after another. Each prisoner may
  open and look into 50 drawers in any order. The drawers are closed again
  afterwards. If, during this search, every prisoner finds his number in one of
  the drawers, all prisoners are pardoned. If some prisoner does not find his
  number, all prisoners die. Before the first prisoner enters the room, the
  prisoners may discuss strategy, but they cannot communicate once the first
  prisoner enters the room.
  Which of the following strategies is better for the prisoners? 
  \begin{enumerate}
  \item Strategy 1 (simple): every prisoner selects 50 drawers at random.
  \item Strategy 2 (complicated): 
    \begin{enumerate}
    \item Each prisoner first opens the drawer with his own number.
    \item If this drawer contains his number he is done and was
      successful.
    \item Otherwise, the drawer contains the number of another prisoner
      and he next opens the drawer with this number.
    \item The prisoner repeats steps 2 and 3 until he finds his own number
      or has opened 50 drawers.
    \end{enumerate}
  \end{enumerate}
\end{enumerate}

% Other possible topics:
% \begin{itemize}
% \item Poker Probabilities: \url{https://en.wikipedia.org/wiki/Poker_probability}
% \item Blackjack Probability: \url{https://en.wikipedia.org/wiki/Blackjack}
% \end{itemize}


\section{Runs and Wins in Baseball}
The percentage of wins obtained by a team over the course of a season
is strongly related with the number of runs it scores and allows. The
relationship between runs and wins is explored in Chapter~4 of the
book by
\href{https://www.amazon.com/gp/product/B07KRNP2BB/ref=dbs_a_def_rwt_bibl_vppi_i0}{Marchi,
  Albert, and Baumer (2019)}.
The source code of the book is at
\url{https://github.com/bayesball/bayesball.github.io}.
A workshop on baseball analytics in a past UConn Sports
Analytics Symposium (UCSAS) led by Dr. Zhe Wang included this as an
example. The training materials of this workshop can be accessed at
\url{https://statds.org/events/ucsas2020/workshops.html#baseball}.
In particular, after downloading it, the code for this example is
\texttt{Runs\_and\_Wins.R} under the \texttt{Code} folder.

\begin{enumerate}
\item
Run the example code on your own computer line by line to reproduce
the analysis.

\item
Explain what Pythagorean formula is and how well it fits baseball.

\item
Follow your interests to further explore the data and report your
findings. Examples topics are: relationship between win percentage and
run differential across decades; Pythagorean residuals for stronger
and weaker teams. Be creative and don't be limited by the
examples.
\end{enumerate}


\section{Crashes in NYC}
The NYC motor vehicle collisions crash table contains details on the crashes.
The real-time full data with documentation is available from 
\href{https://data.cityofnewyork.us/Public-Safety/Motor-Vehicle-Collisions-Crashes/h9gi-nx95}{NYC
  Open Data}.
The dataset that we consider here contains only the crashes in
\href{https://github.com/statds/ids-s23/blob/main/data/nyc_crashes_202301_cleaned.csv}{January,
  2023}, in CSV format, which has been cleaned to some extent.
You can read into R with function \texttt{read.csv()} with the URL as the input.

\begin{enumerate}
\item
  Construct a contigency table for missing in geocode (latitude and longitude)
  by borough. Is the missing pattern the same across borough? Perform some
  hypothesis test if you can.
\item
  Construct an \texttt{hour} variable with integer values from 0 to 23. Plot the
  histogram of the number of crashes by hour. Plot it by borough.
\item
  Create a new variable \texttt{injury} which is one if the number of persons
  injured is~1 or more; and zero otherwise. Construct a cross table for injury
  versus borough. Do the patter look similar across boroughs? Devise a
  hypothesis and test it.
\item
  Contributing factor vehicle~1 contains the first factor of the collision. Get
  a frequency table of the contributing factors.
\item
  Take a subset of the dataset that contains only crashes due to the top~3
  contributing factors. Construct a contingency table for \texttt{injury} by
  the top~3 contributing factors. Device a test to check if injury is associated
  with the contributing factors.  
\end{enumerate}


\end{document}
